\documentclass[10pt,a4paper]{amsart}
\usepackage[margin=19mm]{geometry}
\usepackage{amsmath,amssymb,mathtools}
\usepackage[scr=rsfs,cal=euler]{mathalfa}
\usepackage{xfrac}
\usepackage[unicode,hidelinks,bookmarks=false]{hyperref}
\newcommand{\ie}{i.e.\ }
\newcommand{\cf}{cf.\ }
\newcommand{\resp}{respectively\ }
\newcommand{\Subsection}{\subsection}
\providecommand{\nobreakdash}{\nobreak\hbox{-}}
\setlength{\emergencystretch}{2em}
\multlinegap0pt
\usepackage{bm}
\usepackage{bbm}
\usepackage{dsfont}
\usepackage{xspace}
\usepackage{cancel}
\usepackage{mathtools}
\usepackage{enumitem}
\usepackage{amsthm}
\makeatletter
\@ifundefined{thm}{\newtheorem{thm}{Thm}}{}
\@ifundefined{lem}{\newtheorem{lem}[thm]{Lemma}}{}
\makeatother
\title[Ordinary Isogeny Graphs with Level Structure]{Ordinary Isogeny Graphs with Level Structure}
\author{Derek Perrin, Jos\'e Felipe Voloch}
\date{Source: arXiv:2411.02732v2 (CC BY 4.0)}
\begin{document}
\maketitle

\noindent\emph{Source and licence.} Ordinary Isogeny Graphs with Level Structure. Version arXiv:2411.02732v2 (CC BY 4.0), distributed under the Creative Commons Attribution 4.0 International licence, as recorded in the arXiv licence metadata for this exact version and shown on the abstract page https://arxiv.org/abs/2411.02732v2. Full rights notice: licenses/arxiv-2411.02732-CC-BY-4.0.txt.\par\smallskip

\noindent\textit{Editorial scope.} Counted unit: the Lem whose source label is \texttt{lem:volcano-growth} in the article (source0.tex lines 204--224, source label \texttt{lem:volcano-growth}), delivered together with the article's own complete original proof of it (source0.tex lines 225--253, one \texttt{proof} environment of 29 lines). No mathematics is written, repaired or re-derived: statement and proof are the article's own text line for line. Required context retained uncounted: none. Every definition, display and label of those lines is retained unchanged. Omitted from this package and not redistributed: 1--203, 254--707. Nothing inside a retained excerpt is omitted or altered: every retained body is delivered exactly as the article's lines are written, so no omission receipt of any kind accompanies this package. Citations: the retained excerpts carry no citation command at all, so no bibliography entry of the article is transcribed and no reference is redistributed. Reference adaptation: none was needed and none was made; every reference inside the delivered unit resolves inside it, so no unresolved reference or citation is produced. Printed slips kept literal: no printed word, symbol, hypothesis or number of the counted statement or of its proof was altered. Layout and class: the article's own class is \texttt{amsart} with packages including fontenc, setspace, hyperref, mathtools, amsthm, amssymb, cleveref, enumitem, xcolor, todonotes, tikz, ifthen; the wrapper is the lane's pinned \texttt{amsart} editorial wrapper at 10pt on A4, which restates the theorem-like environments the retained text needs and adds the article's own macro definitions that the retained text needs (none), with extra wrapper packages (bm, bbm, dsfont, xspace, cancel, mathtools, enumitem). The delivered numbering is the unit's own. Assets: no retained span contains a figure environment, a graphics-inclusion command, a PDF-page inclusion, a table or a TikZ picture; no image file is copied into this package, the package declares no asset dependency and no image file is redistributed with it. Licence: version arXiv:2411.02732v2 of this article is distributed under the Creative Commons Attribution 4.0 International licence, as recorded in the arXiv licence metadata for this exact version and shown on the abstract page https://arxiv.org/abs/2411.02732v2. A full rights notice is delivered at licenses/arxiv-2411.02732-CC-BY-4.0.txt. Characters: every retained excerpt is pure ASCII, so the delivered engine, XeLaTeX with its default Latin Modern text font, reports no missing character.
\begin{lem}\label{lem:volcano-growth}
  Let $G_{\ell}$ be an isogeny volcano with $n$ vertices.
  \begin{enumerate}[label=(\roman*)]
  \item The graph $G_{\ell,0}(N)$ has
    \[
      n N\prod_{p \mid N}\left( 1 + \frac{1}{p} \right)
    \]
    vertices.
  \item If $N > 2,$ the graph $G_{\ell,1}(N)$ has
    \[
      \frac{nN\phi(N)}{2}\prod_{p \mid N}\left( 1 + \frac{1}{p} \right)
    \]
    vertices and $3n$ if $N = 2$.
  \item If $N > 2,$ the graph $G_{\ell}(N)$ has
    \[
      \frac{nN^2\phi(N)^2}{2} \prod_{p \mid N}\left( 1 + \frac{1}{p} \right)
    \]
    vertices and $6n$ if $N = 2$.
  \end{enumerate}
  where $\phi$ is Euler's totient function.
\end{lem}

\begin{proof}
  For each case, we proceed by counting the number of level structures which can be added to each vertex in $G_{\ell}.$

  We begin by proving (ii). Since $E[N] \simeq \mathbb{Z}/N\mathbb{Z} \times \mathbb{Z}/N \mathbb{Z}$ and $N = \prod p_i^{e_i},$ then by the Chinese Remainder Theorem we may write
  \begin{equation}\label{eq:torsion-factors}
    E[N] \simeq \left( \mathbb{Z}/p_1^{e_1}\mathbb{Z} \right)^2 \times \left( \mathbb{Z}/p_2^{e_2}\mathbb{Z} \right)^2 \times \cdots \times \left( \mathbb{Z}/p_r^{e_r}\mathbb{Z} \right)^2.
  \end{equation}
   An element $P$ of $E[N]$ has maximal order $N$ if and only if it has maximal order in each of factors on the right-hand side of \eqref{eq:torsion-factors}, and so it is sufficient to find the number of elements of maximal order of each factor of $E[N].$ An element $(a,b)$ of $(\mathbb{Z}/p_i^{e_i})^2$ is of maximal order if either $(a,p_i^{e_i}) = 1$ or $(b, p_i^{e_i}) = 1.$ There are $\phi(p_i^{e_i})p_i^{e_i}$ pairs where $a$ is coprime to $p_i^{e_i}$ and $\phi(p_i^{e_i})(p_i^{e_i} - \phi(p_i^{e_i}))$ pairs where $b$ is coprime to $p_i^{e_i}$ and $a$ is not. Some algebraic manipulation shows there are a total of $p_i^{2e_i} - p_i^{2(e_i-1)}$ such elements of maximal order in $(\mathbb{Z}/p_i^{e_i}\mathbb{Z})^2$. We repeat this for each prime power factor of $N$ and by the Chinese Remainder Theorem we may take their product to get
  \begin{equation}\label{eq:num-points}
    N \phi(N) \prod_{p \mid N}\left(1 + \frac{1}{p}\right)
  \end{equation}
  total points of order $N.$ If $N = 2,$ then $P = -P$ for any point $P$ and so we may substitute $N = 2$ into \eqref{eq:num-points} to get a scaling factor of $3.$ Otherwise, if $N > 2,$ then we must divide the result of \eqref{eq:num-points} by $2$ due to identifying $P$ with $-P.$

  To prove (i), we use the fact that any cyclic subgroup $\left\langle P \right\rangle \subseteq E[N]$ of order $N$ has $\phi(N)$ elements with order $N.$ We may therefore partition the elements $P$ of order $N$ into equivalence classes where $P \sim P'$ if there exists an integer $k$ such that $P' = kP.$ From (ii), we have that there are
  \[
    N\phi(N)\prod_{p \mid N}\left( 1 + \frac{1}{p} \right)
  \]
  elements of order $N$ in $E[N].$ Dividing by $\phi(N)$ the size of each partition, we see there are
  \[
     N \prod_{p \mid N}\left( 1 + \frac{1}{p} \right)
  \]
  such equivalence classes, and therefore subgroups, of order $N.$

  To prove (iii), we first work with a factor $(\mathbb{Z}/p_i^{e_i}\mathbb{Z})^2$ of \eqref{eq:torsion-factors}. The result of (ii) gives us the number of ways to choose a first basis point $P$. The restrictions on choosing a second basis point $Q$ is that it must be of order $p_i^{e_i}$ and it cannot be congruent to a scalar multiple of $P$ modulo $p_i.$ The congruence requirement arises from the fact that $\mathbb{Z}/p_i^{e_i} \mathbb{Z}$ is local at $p_i.$ From (ii), we know there are $\phi(p_i^{e_i})$ elements of order $p_i^{e_i}$ in $\left\langle P \right\rangle$ and that each of those elements is congruent to $p_i^{e_i-1}$ distinct elements of order $p_i^{e_i}.$ This gives a total of  $\phi(p_i^{e_i})(p_i^{e_i} + p_i^{e_i-1}) - \phi(p_i^{e_i})p_i^{e_i-1} = p_i^{e_i}\phi(p_i^{e_i})$ choices for $Q.$ Taking the product over all $p_i$ by the Chinese Remainder Theorem gives
  \begin{equation}\label{eq:num-bases}
    N^2\phi(N)^2\prod_{p \mid N}\left( 1 + \frac{1}{p} \right)
  \end{equation}
  choices of bases. If $N = 2,$ then $P = -P$ and $Q = -Q$ for any choice of $P, Q,$ and so substituting $N = 2$ into \eqref{eq:num-bases} gives a scaling factor of $6.$ If $N > 2,$ we divide \eqref{eq:num-bases} by $2$ due to identifying the pair $(P, Q)$ with $(-P, -Q).$
\end{proof}

\end{document}
